\section{Related Work}
\label{sec:related}

\para{Pragmatic competence in LLMs.}
A growing body of work evaluates whether LLMs understand pragmatic phenomena. \citet{ruis2024goldilocks} introduce the LUDWIG benchmark for conversational implicature and show that GPT-4 achieves 88.7\% accuracy with chain-of-thought prompting, while smaller models perform near chance. \citet{cho2024pragmatic} probe BERT and GPT-2 for scalar implicature inference and find that both models prefer pragmatic interpretations in neutral contexts, with GPT-2 showing increased surprisal when pragmatic inference is required under upper-bound questions. \citet{lin2024probing} compare direct (embedding-based) and indirect (perplexity-based) methods for probing scalar adjective knowledge, finding that LLMs encode rich lexical-semantic information but show ``unsatisfying performance on scalar diversity''---the pragmatic aspect. \citet{sato2025pragmatic} demonstrate that providing pragmatic theory descriptions (\eg Grice's maxims) as prompts improves LLM performance on implied-meaning tasks by up to 9.6\%, suggesting that models possess latent pragmatic knowledge that is not always activated. Collectively, these studies establish that LLMs \emph{can} reason about pragmatic phenomena but leave open whether they \emph{do so} during unconstrained generation.

\para{Perplexity and surprisal as probes.}
Using model-internal probability estimates to probe linguistic knowledge is well established. \citet{lin2024probing} compare perplexity of correct vs.\ incorrect scalar constructions (\eg ``good but not awesome'' vs.\ ``awesome but not good'') to assess whether models encode intensity orderings. \citet{cho2024pragmatic} measure surprisal under different question contexts to detect pragmatic processing. \citet{cifka2023blackbox} propose context length probing, which tracks how a model's predictions change as context is incrementally extended, yielding importance scores via KL divergence. We build on this tradition but shift from probing \emph{comprehension} to probing \emph{generation preferences}: rather than asking whether the model assigns lower perplexity to linguistically correct constructions, we ask whether it shows systematic directional preferences along pragmatic axes.

\para{Behavioral vs.\ introspective assessment.}
\citet{trienes2025behavioral} use length-constrained summarization as a behavioral probe for information salience, tracking which content an LLM retains at different compression levels. They find that models' behavioral salience hierarchies are consistent across model families but do \emph{not} match their self-reported salience when asked directly. This behavioral-introspective gap motivates our comparison of perplexity-detected pragmatic preferences against introspective variable reports. Our work extends their finding from content salience to pragmatic variable consideration, using a complementary probing method (log-likelihood of pre-filled responses rather than constrained generation).

\para{Questions Under Discussion.}
The QUD framework from formal pragmatics \citep{roberts2012information} provides a theoretical grounding for what is ``under consideration'' in discourse. Recent computational work has applied QUD parsing to NLP tasks \citep{wu2024qudeval, ko2024qudselect}, and surveys have connected QUD models to language generation \citep{westera2025survey}. While we do not directly implement a QUD parser, our notion of ``variables under consideration'' is conceptually aligned: a pragmatic variable is under consideration if the model's generation prior is sensitive to variation along that dimension, analogous to how a QUD constrains what information is relevant in a discourse.
